\chapter{Conclusiones}
\label{ch:conclusiones}

Los resultados no apuntan a un único controlador ganador para toda la planta. En el lazo servo de Fase~1, las dos políticas aprendidas reducen casi a la mitad el IAE del PID (QL $1.80$ y DL $1.90$ frente a $3.89$), mientras que el difuso queda por detrás ($5.97$) al no incluir acción integral. Esa lectura, sin embargo, cambia al mirar la señal de mando: el QL alcanza el seguimiento con una acción bang-bang que satura el actuador y acumula la mayor variación del grupo; la red profunda mantiene un error parecido con la señal más suave y sin saturación. Para un actuador físico con límite de velocidad de cambio, esa diferencia pesa más que la décima de IAE y justifica escoger la red en el servo local.

En la celda completa conviene decidir con más de una métrica. El PID supervisado, una vez resintonizado, consigue el mejor seguimiento global (IAE $11.4$) y el menor desbalance entre esteras. Las políticas aprendidas no lo igualan en error, pero sí aportan perfiles útiles: el DL prioriza margen de seguridad y productividad, y el QL reduce consumo y apenas se degrada en el escenario severo. El difuso vuelve a perder en IAE por la falta de integral, aunque conserva dos ventajas prácticas para una primera puesta en marcha: mantiene el atasco bajo control y deja una base de reglas entendible para el personal de planta.

La robustez fue el filtro más exigente del trabajo. Sin la capa supervisora, el PID desborda la celda en el escenario severo (3.64 piezas y 20 muestras en alarma); con ella, ninguno de los cuatro métodos viola la capacidad en los ensayos evaluados. La rampa trapezoidal también dejó una advertencia concreta para el aprendizaje por refuerzo tabular: representar la perturbación como variable binaria crea un punto ciego de seis muestras y hace que el QL, pese a su buen resultado servo, sea el único método que empeora ante una perturbación gradual. En este banco, la codificación del estado resultó tan importante como el algoritmo elegido.

La propuesta final queda dividida por nivel de control. Para el servo local se recomienda la red profunda entrenada sobre el gemelo digital, porque combina menor error que el PID, mando suave, ausencia de saturación e inferencia de coste fijo por ciclo. Para la primera implantación en la celda completa se recomienda el controlador difuso junto con la capa supervisora de capacidad: sacrifica parte del IAE, pero ofrece trazabilidad, ajuste directo por personal de planta y restricciones de seguridad explícitas. El PID-Aware queda como referencia industrial y respaldo operativo; DL y QL deberían ejecutarse inicialmente en modo sombra sobre el gemelo digital antes de autorizar su uso en línea.

El banco desarrollado también queda como resultado reutilizable. Al comparar todos los controladores con la misma planta, los mismos escenarios y las mismas semillas, las diferencias observadas se pueden atribuir al controlador y no al experimento. Esa trazabilidad permitió comprobar que ampliar la red y el entrenamiento de Fase~2 redujo el IAE de forma sustancial, señal de que el límite estaba en la capacidad del aproximador y no en el enfoque. La misma infraestructura servirá para evaluar las extensiones del Capítulo~\ref{ch:i40} antes de intervenir la línea real.
